%!TEX root =  tfg.tex
\chapter{Objetivos}

\begin{quotation}[Oceanographer]{Jacques Yves Cousteau}
The sea, once it casts its spell, holds one in its net of wonder forever.
\end{quotation}

\begin{abstract}
Visto el contexto, formalizaremos los objetivos de la app de forma concreta. ¿Hacía quien esta dirigida la aplicación? ¿Qué funcionalidades tendrá? ¿Qué tecnologías se van a usar? ¿Qué se va a conseguir con el proyecto?
\end{abstract}

\section{Motivación}
Como ya hemos comentado, la idea de Scubex es centralizar funcionalidades para el nicho del buceo. Para lograrlo se han planteado esta serie de objetivos que, una vez cumplidos, darán como resultado la aplicación deseada.

\section{Listado de objetivos}

\begin{description}
\item \textbf{Objetivo 1. Implementar un sistema de autenticación de usuarios} \\
Desarrollar un sistema de autenticación seguro, permitiendo a los usuarios registrarse e iniciar sesión de forma sencilla. El sistema debe guardar los perfiles de usuario y sus datos (nombre, foto, incursiones, comentarios).

\item \textbf{Objetivo 2. Crear un mapa interactivo de exploración marina} \\
Integrar un mapa interactivo con controles de zoom y navegación que permita a los usuarios explorar zonas de buceo de forma intuitiva y fluida.

\item \textbf{Objetivo 3. Integrar datos científicos de biodiversidad marina} \\
Desarrollar un escáner de especies que consulte APIs científicas para identificar fauna marina en una zona determinada. 

\item \textbf{Objetivo 4. Proporcionar información climatológica y oceanográfica en tiempo real} \\
Ofrecer datos actualizados sobre condiciones climatológicas de la zona: viento, temperatura del agua, corrientes marinas, altura de olas y datos atmosféricos relevantes para la planificación de inmersiones en un punto de buceo concreto.

\item \textbf{Objetivo 5. Implementar un sistema de registro de publicaciones} \\
Permitir a los usuarios publicar sus inmersiones en el mapa interactivo: título, descripción opcional, fotografía y coordenadas geográficas. La localización se enriquece automáticamente con el nombre del lugar (playa, municipio, estado) mediante geocodificación inversa a través de la API Nominatim de OpenStreetMap, con una estrategia de cascada de tres niveles de zoom para garantizar la obtención del nombre más relevante. El autor puede editar y eliminar sus propias publicaciones.

\item \textbf{Objetivo 6. Crear un sistema de interacción social sobre publicaciones} \\
Desarrollar un sistema completo de interacciones sobre publicaciones: \textbf{(1)} likes (con contador visible), \textbf{(2)} guardado de publicaciones para acceso posterior desde el perfil, y \textbf{(3)} comentarios con soporte de menciones en formato \texttt{@[Nombre](email)}. Las menciones generan notificaciones automáticas al usuario mencionado. Cada usuario puede eliminar sus propios comentarios.

\item \textbf{Objetivo 7. Implementar una red social de seguimiento entre usuarios} \\
Desarrollar un sistema de perfiles públicos y seguimiento entre usuarios: cada perfil muestra nombre, foto, número de seguidores y seguidos, y grid de publicaciones. Cualquier usuario puede seguir o dejar de seguir a otro. El perfil propio incluye pestañas de publicaciones propias y publicaciones guardadas. Los perfiles son navegables directamente desde el mapa (autor de una publicación) o desde comentarios (autor de un comentario).

\item \textbf{Objetivo 8. Implementar un sistema de notificaciones en tiempo real} \\
Proporcionar retroalimentación inmediata al usuario sobre eventos relevantes: nuevos seguidores, likes recibidos, comentarios en sus publicaciones y menciones en comentarios ajenos. Las notificaciones se muestran mediante una campana en el header con badge de contador, con soporte de lectura individual, marcado masivo como leído y eliminación permanente (swipe en móvil, botón en escritorio).

\end{description}